% ============================================================================
% CAPÍTULO 13: SYSTEMS BIOLOGY IN MEDICINE
% Bioinformática para Biologia de Sistemas
% ============================================================================

% ============================================================================
% TEMPLATE BASE PARA APRESENTAÇÕES EM BEAMER
% Bioinformática para Biologia de Sistemas
% ============================================================================

\documentclass[12pt, aspectratio=169]{beamer}

% ============================================================================
% PACOTES ESSENCIAIS
% ============================================================================
\usepackage[utf8]{inputenc}
\usepackage[T1]{fontenc}
\usepackage[brazilian]{babel}
\usepackage{amsmath, amssymb, amsthm}
\usepackage{graphicx}
\usepackage{xcolor}
\usepackage{tikz}
\usetikzlibrary{arrows.meta, shapes, positioning, calc, decorations.pathreplacing, decorations.shapes, decorations.pathmorphing}
\usepackage{listings}
\usepackage{booktabs}
\usepackage{multirow}
\usepackage{hyperref}
\usepackage{chemfig}
\usepackage{chemarr}

% ============================================================================
% CONFIGURAÇÃO DO TEMA
% ============================================================================
\usetheme{Madrid}
\usecolortheme{default}

% Cores personalizadas
\definecolor{azulescuro}{RGB}{0, 51, 102}
\definecolor{azulclaro}{RGB}{51, 153, 255}
\definecolor{verdeacido}{RGB}{102, 204, 0}
\definecolor{laranjacel}{RGB}{255, 153, 51}
\definecolor{roxodna}{RGB}{153, 51, 255}

\setbeamercolor{structure}{fg=azulescuro}
\setbeamercolor{palette primary}{bg=azulescuro,fg=white}
\setbeamercolor{palette secondary}{bg=azulclaro,fg=white}
\setbeamercolor{palette tertiary}{bg=azulescuro,fg=white}
\setbeamercolor{palette quaternary}{bg=azulclaro,fg=white}
\setbeamercolor{block title}{bg=azulescuro,fg=white}
\setbeamercolor{block body}{bg=azulclaro!10,fg=black}
\setbeamercolor{block title alerted}{bg=red!80,fg=white}
\setbeamercolor{block body alerted}{bg=red!10,fg=black}
\setbeamercolor{block title example}{bg=verdeacido!80,fg=white}
\setbeamercolor{block body example}{bg=verdeacido!10,fg=black}

% ============================================================================
% CONFIGURAÇÃO DE CÓDIGO
% ============================================================================
\lstset{
    language=Python,
    basicstyle=\ttfamily\small,
    keywordstyle=\color{azulescuro}\bfseries,
    commentstyle=\color{verdeacido},
    stringstyle=\color{laranjacel},
    numbers=left,
    numberstyle=\tiny\color{gray},
    stepnumber=1,
    numbersep=5pt,
    backgroundcolor=\color{white},
    showspaces=false,
    showstringspaces=false,
    showtabs=false,
    frame=single,
    rulecolor=\color{black},
    tabsize=2,
    captionpos=b,
    breaklines=true,
    breakatwhitespace=false,
    escapeinside={\%*}{*)},
    xleftmargin=10pt,
    xrightmargin=10pt
}

% Definir estilo para R
\lstdefinestyle{Rstyle}{
    language=R,
    basicstyle=\ttfamily\small,
    keywordstyle=\color{azulescuro}\bfseries,
    commentstyle=\color{verdeacido},
    stringstyle=\color{laranjacel}
}

% Definir estilo para MATLAB
\lstdefinestyle{MATLABstyle}{
    language=Matlab,
    basicstyle=\ttfamily\small,
    keywordstyle=\color{azulescuro}\bfseries,
    commentstyle=\color{verdeacido},
    stringstyle=\color{laranjacel}
}

% ============================================================================
% COMANDOS PERSONALIZADOS
% ============================================================================

% Comando para equações destacadas
\newcommand{\destaque}[1]{%
    \begin{alertblock}{}
        \begin{center}
            #1
        \end{center}
    \end{alertblock}
}

% Comando para definições
\newcommand{\definicao}[2]{%
    \begin{block}{Definição: #1}
        #2
    \end{block}
}

% Comando para exemplos
\newcommand{\exemplo}[2]{%
    \begin{exampleblock}{Exemplo: #1}
        #2
    \end{exampleblock}
}

% Comando para observações importantes
\newcommand{\importante}[1]{%
    \begin{alertblock}{Importante!}
        #1
    \end{alertblock}
}

% Comandos matemáticos úteis
\newcommand{\R}{\mathbb{R}}
\newcommand{\N}{\mathbb{N}}
\newcommand{\Z}{\mathbb{Z}}
\newcommand{\dif}{\mathrm{d}}
\newcommand{\parcial}[2]{\frac{\partial #1}{\partial #2}}
\newcommand{\derivada}[2]{\frac{\dif #1}{\dif #2}}
\newcommand{\vect}[1]{\boldsymbol{#1}}

% Comandos biológicos
\newcommand{\gene}[1]{\textit{#1}}
\newcommand{\proteina}[1]{\textsc{#1}}
\newcommand{\especie}[1]{\textit{#1}}

% ============================================================================
% CONFIGURAÇÃO DE RODAPÉ
% ============================================================================
\setbeamertemplate{footline}{
  \leavevmode%
  \hbox{%
  \begin{beamercolorbox}[wd=.33\paperwidth,ht=2.25ex,dp=1ex,center]{author in head/foot}%
    \usebeamerfont{author in head/foot}\insertshortauthor
  \end{beamercolorbox}%
  \begin{beamercolorbox}[wd=.34\paperwidth,ht=2.25ex,dp=1ex,center]{title in head/foot}%
    \usebeamerfont{title in head/foot}\insertshorttitle
  \end{beamercolorbox}%
  \begin{beamercolorbox}[wd=.33\paperwidth,ht=2.25ex,dp=1ex,right]{date in head/foot}%
    \usebeamerfont{date in head/foot}\insertshortdate{}\hspace*{2em}
    \insertframenumber{} / \inserttotalframenumber\hspace*{2ex}
  \end{beamercolorbox}}%
  \vskip0pt%
}

% ============================================================================
% CONFIGURAÇÃO DE NAVEGAÇÃO
% ============================================================================
\setbeamertemplate{navigation symbols}{}

% ============================================================================
% CONFIGURAÇÃO DE ITEMIZE/ENUMERATE
% ============================================================================
\setbeamertemplate{itemize items}[circle]
\setbeamertemplate{enumerate items}[default]

% ============================================================================
% FIM DO TEMPLATE
% ============================================================================


% ============================================================================
% INFORMAÇÕES DO DOCUMENTO
% ============================================================================
\title[Cap. 13: Systems Biology in Medicine]{Capítulo 13: Systems Biology in Medicine}
\subtitle{Bioinformática para Biologia de Sistemas}
\author{Seu Nome}
\institute{Sua Instituição}
\date{\today}

\begin{document}

% ============================================================================
% SLIDE DE TÍTULO
% ============================================================================
\begin{frame}
\titlepage
\end{frame}

% ============================================================================
% SUMÁRIO
% ============================================================================
\begin{frame}{Sumário}
\tableofcontents
\end{frame}

% ============================================================================
% SEÇÃO 1: INTRODUÇÃO
% ============================================================================
\section{Introdução}

\begin{frame}{Visão Geral do Capítulo}
\begin{block}{Objetivos de Aprendizagem}
\begin{itemize}
    \item Compreender medicina de sistemas e medicina personalizada
    \item Aplicar biologia de sistemas à descoberta de fármacos
    \item Modelar doenças complexas como sistemas integrados
    \item Explorar farmacogenômica e estratificação de pacientes
    \item Analisar estudos de caso clínicos com abordagens sistêmicas
\end{itemize}
\end{block}

\vspace{0.3cm}

\begin{exampleblock}{Por que Systems Medicine?}
\begin{itemize}
    \item Doenças complexas requerem abordagens integrativas
    \item Medicina personalizada baseada em perfis moleculares
    \item Descoberta de fármacos mais eficiente e segura
    \item Previsão de respostas terapêuticas individualizadas
\end{itemize}
\end{exampleblock}
\end{frame}

\begin{frame}{Da Bancada ao Leito: Medicina de Sistemas}
\definicao{Medicina de Sistemas}{
Aplicação dos princípios da biologia de sistemas à medicina clínica, integrando dados multi-ômicos, redes moleculares e modelagem computacional para compreender, diagnosticar e tratar doenças.
}

\vspace{0.3cm}

\begin{center}
\begin{tikzpicture}[node distance=2.5cm, auto,
    box/.style={rectangle, draw, fill=azulclaro!20, text width=2.2cm, text centered, rounded corners, minimum height=1cm}]

    \node[box, fill=roxodna!20] (omics) {Dados\\Ômicos};
    \node[box, fill=verdeacido!20, right of=omics] (network) {Redes\\Moleculares};
    \node[box, fill=laranjacel!20, below of=network] (model) {Modelos\\Computacionais};
    \node[box, fill=azulclaro!40, left of=model] (clinical) {Aplicação\\Clínica};

    \draw[->, thick, azulescuro] (omics) -- (network);
    \draw[->, thick, azulescuro] (network) -- (model);
    \draw[->, thick, azulescuro] (model) -- (clinical);
    \draw[->, thick, azulescuro, bend left=45] (clinical) to node[above, font=\tiny] {Validação} (omics);

\end{tikzpicture}
\end{center}
\end{frame}

\begin{frame}{Medicina Personalizada e de Precisão}
\begin{columns}
\column{0.5\textwidth}
\begin{block}{Medicina Tradicional}
\begin{itemize}
    \item Abordagem "one-size-fits-all"
    \item Tratamento baseado em sintomas
    \item Decisões empíricas
    \item Resposta variável
\end{itemize}
\end{block}

\column{0.5\textwidth}
\begin{block}{Medicina Personalizada}
\begin{itemize}
    \item Tratamento individualizado
    \item Baseada em perfil molecular
    \item Decisões orientadas por dados
    \item Resposta otimizada
\end{itemize}
\end{block}
\end{columns}

\vspace{0.5cm}

\importante{
Medicina personalizada = tratamento certo + paciente certo + momento certo
}

\begin{exampleblock}{Pilares da Medicina de Precisão}
\begin{itemize}
    \item Genômica individual
    \item Biomarcadores moleculares
    \item Estratificação de pacientes
    \item Farmacogenômica
\end{itemize}
\end{exampleblock}
\end{frame}

\begin{frame}{Abordagens Sistêmicas para Doenças}
\begin{block}{Mudança de Paradigma}
\begin{center}
\begin{tabular}{lcc}
\toprule
\textbf{Aspecto} & \textbf{Tradicional} & \textbf{Sistêmica} \\
\midrule
Foco & Gene/proteína única & Redes moleculares \\
Causa & Linear & Multifatorial \\
Terapia & Alvo único & Multi-alvo \\
Dados & Isolados & Integrados \\
Modelo & Reducionista & Holístico \\
\bottomrule
\end{tabular}
\end{center}
\end{block}

\vspace{0.3cm}

\begin{exampleblock}{Aplicações Clínicas}
\begin{itemize}
    \item Identificação de alvos terapêuticos
    \item Reposicionamento de fármacos
    \item Biomarcadores diagnósticos e prognósticos
    \item Estratificação de pacientes em ensaios clínicos
    \item Previsão de efeitos adversos
\end{itemize}
\end{exampleblock}
\end{frame}

% ============================================================================
% SEÇÃO 2: DESCOBERTA E DESENVOLVIMENTO DE FÁRMACOS
% ============================================================================
\section{Descoberta e Desenvolvimento de Fármacos}

\begin{frame}{Descoberta de Fármacos: Tradicional vs Sistêmica}
\begin{columns}
\column{0.5\textwidth}
\textbf{Abordagem Tradicional:}
\begin{enumerate}
    \item Identificar alvo único
    \item Triagem de compostos
    \item Otimização de lead
    \item Testes pré-clínicos
    \item Ensaios clínicos
\end{enumerate}

\vspace{0.3cm}

\textbf{Limitações:}
\begin{itemize}
    \item Alta taxa de falha (>90\%)
    \item Custo elevado (\$2.6 bilhões)
    \item Tempo longo (10-15 anos)
    \item Efeitos adversos inesperados
\end{itemize}

\column{0.5\textwidth}
\textbf{Abordagem Sistêmica:}
\begin{enumerate}
    \item Análise de redes
    \item Identificação de módulos
    \item Alvos múltiplos/sinérgicos
    \item Modelagem \textit{in silico}
    \item Estratificação de pacientes
\end{enumerate}

\vspace{0.3cm}

\textbf{Vantagens:}
\begin{itemize}
    \item Melhor compreensão de mecanismos
    \item Previsão de efeitos adversos
    \item Reposicionamento de fármacos
    \item Medicina personalizada
\end{itemize}
\end{columns}
\end{frame}

\begin{frame}{Identificação e Validação de Alvos}
\begin{block}{Pipeline Sistêmico}
\begin{enumerate}
    \item \textbf{Análise de redes de doenças}
    \begin{itemize}
        \item Identificar genes/proteínas associados à doença
        \item Análise de centralidade (degree, betweenness, closeness)
        \item Identificar hubs e bottlenecks
    \end{itemize}

    \item \textbf{Módulos de doença}
    \begin{itemize}
        \item Detecção de comunidades
        \item Módulos funcionalmente relacionados
        \item Proximidade de rede
    \end{itemize}

    \item \textbf{Critérios de validação}
    \begin{itemize}
        \item Druggability (capacidade de ser alvo de fármaco)
        \item Especificidade tecidual
        \item Papel funcional validado
        \item Sem efeitos adversos graves
    \end{itemize}
\end{enumerate}
\end{block}
\end{frame}

\begin{frame}{Farmacologia de Redes (Network Pharmacology)}
\definicao{Network Pharmacology}{
Estudo de interações fármaco-alvo no contexto de redes biológicas, reconhecendo que fármacos afetam múltiplos alvos e vias.
}

\vspace{0.3cm}

\begin{center}
\begin{tikzpicture}[scale=0.8]
    % Proteinas alvo
    \node[draw, circle, fill=verdeacido!30, minimum size=0.8cm] (t1) at (0,0) {T1};
    \node[draw, circle, fill=verdeacido!30, minimum size=0.8cm] (t2) at (2,1) {T2};
    \node[draw, circle, fill=verdeacido!30, minimum size=0.8cm] (t3) at (2,-1) {T3};
    \node[draw, circle, fill=azulclaro!30, minimum size=0.8cm] (t4) at (4,0.5) {T4};
    \node[draw, circle, fill=azulclaro!30, minimum size=0.8cm] (t5) at (4,-0.5) {T5};

    % Farmaco
    \node[draw, diamond, fill=laranjacel!50, minimum size=1cm] (drug) at (-2,0) {\textbf{D}};

    % Interacoes farmaco-alvo
    \draw[->, thick, red] (drug) -- (t1);
    \draw[->, thick, red] (drug) to[bend left=15] (t2);
    \draw[->, thick, red] (drug) to[bend right=15] (t3);

    % Rede proteina-proteina
    \draw[thick, azulescuro] (t1) -- (t2);
    \draw[thick, azulescuro] (t1) -- (t3);
    \draw[thick, azulescuro] (t2) -- (t4);
    \draw[thick, azulescuro] (t3) -- (t5);
    \draw[thick, azulescuro] (t4) -- (t5);

    % Legenda
    \node[below, font=\tiny] at (1,-2) {Alvos primários};
    \node[below, font=\tiny] at (4,-1.5) {Alvos secundários};
\end{tikzpicture}
\end{center}

\importante{
Polifarmacologia: um fármaco, múltiplos alvos $\rightarrow$ efeitos sinérgicos
}
\end{frame}

\begin{frame}{Fármacos Multi-Alvo}
\begin{block}{Conceito}
Fármacos projetados para modular múltiplos alvos simultaneamente, explorando sinergias e redundâncias em redes biológicas.
\end{block}

\begin{columns}
\column{0.5\textwidth}
\textbf{Vantagens:}
\begin{itemize}
    \item Maior eficácia
    \item Menor resistência
    \item Efeitos sinérgicos
    \item Adequado para doenças complexas
\end{itemize}

\vspace{0.3cm}

\textbf{Exemplos:}
\begin{itemize}
    \item Inibidores multi-quinase (câncer)
    \item Aspirina (COX-1/COX-2)
    \item Estatinas (HMG-CoA + pleiotrópicos)
\end{itemize}

\column{0.5\textwidth}
\textbf{Desafios:}
\begin{itemize}
    \item Design complexo
    \item Otimização de múltiplos alvos
    \item Previsão de efeitos adversos
    \item Propriedades farmacocinéticas
\end{itemize}

\vspace{0.3cm}

\begin{exampleblock}{Índice de Seletividade}
\begin{equation*}
S = \frac{IC_{50}(\text{off-target})}{IC_{50}(\text{on-target})}
\end{equation*}

$S > 100$: seletivo\\
$S < 10$: multi-alvo
\end{exampleblock}
\end{columns}
\end{frame}

\begin{frame}{Redes de Interação Fármaco-Alvo}
\begin{block}{Construção de Redes DTI}
\textbf{Dados:}
\begin{itemize}
    \item Estrutura química de fármacos
    \item Estrutura 3D de proteínas alvo
    \item Dados de afinidade (IC$_{50}$, K$_d$)
    \item Screening em larga escala
    \item Bases de dados: DrugBank, ChEMBL, BindingDB
\end{itemize}
\end{block}

\vspace{0.3cm}

\begin{columns}
\column{0.5\textwidth}
\textbf{Análises:}
\begin{itemize}
    \item Grau de promiscuidade de fármacos
    \item Polífarmacologia
    \item Reposicionamento
    \item Previsão de efeitos adversos
\end{itemize}

\column{0.5\textwidth}
\begin{exampleblock}{Métrica de Promiscuidade}
\begin{equation*}
P_d = \frac{\text{\# alvos do fármaco}}{\text{\# total de fármacos}}
\end{equation*}

Fármacos com $P_d$ alto:
\begin{itemize}
    \item Maior chance de efeitos adversos
    \item Potencial para reposicionamento
\end{itemize}
\end{exampleblock}
\end{columns}
\end{frame}

\begin{frame}{Farmacocinética: Modelo ADME}
\begin{block}{ADME}
\begin{itemize}
    \item \textbf{A}bsorção: entrada do fármaco no organismo
    \item \textbf{D}istribuição: dispersão pelos tecidos
    \item \textbf{M}etabolismo: biotransformação (fígado)
    \item \textbf{E}xcreção: eliminação (rins, bile)
\end{itemize}
\end{block}

\vspace{0.3cm}

\begin{columns}
\column{0.5\textwidth}
\textbf{Fatores que afetam ADME:}
\begin{itemize}
    \item Propriedades físico-químicas
    \item Solubilidade, permeabilidade
    \item Ligação a proteínas plasmáticas
    \item Metabolismo de primeira passagem
    \item Polimorfismos genéticos (CYP450)
\end{itemize}

\column{0.5\textwidth}
\begin{alertblock}{Regra de Lipinski (Regra dos 5)}
Fármaco oralmente biodisponível:
\begin{itemize}
    \item Peso molecular < 500 Da
    \item LogP < 5
    \item Doadores H-bond $\leq$ 5
    \item Aceptores H-bond $\leq$ 10
\end{itemize}
\end{alertblock}
\end{columns}
\end{frame}

\begin{frame}{Modelagem PK/PD}
\definicao{Farmacocinética/Farmacodinâmica}{
\textbf{PK:} "O que o corpo faz com o fármaco" --- concentração vs tempo\\
\textbf{PD:} "O que o fármaco faz com o corpo" --- efeito vs concentração
}

\vspace{0.3cm}

\begin{block}{Modelos Compartimentais}
Representam o corpo como compartimentos onde o fármaco se distribui.

\textbf{Tipos:}
\begin{itemize}
    \item Modelo de 1 compartimento (mais simples)
    \item Modelo de 2 compartimentos (central + periférico)
    \item Modelo de 3+ compartimentos (tecidos específicos)
\end{itemize}
\end{block}

\begin{exampleblock}{Parâmetros PK Importantes}
\begin{itemize}
    \item $C_{max}$: concentração máxima
    \item $T_{max}$: tempo de concentração máxima
    \item AUC: área sob a curva (exposição total)
    \item $t_{1/2}$: meia-vida de eliminação
    \item $V_d$: volume de distribuição
    \item CL: clearance (depuração)
\end{itemize}
\end{exampleblock}
\end{frame}

\begin{frame}{Modelo de Um Compartimento}
\begin{columns}
\column{0.5\textwidth}
\begin{center}
\begin{tikzpicture}[scale=0.9]
    % Compartimento central
    \draw[thick, azulescuro, fill=azulclaro!20] (0,0) rectangle (3,2);
    \node at (1.5,1) {\textbf{Compartimento}};
    \node at (1.5,0.5) {\textbf{Central}};
    \node at (1.5,1.7) {$V_d$, $C(t)$};

    % Entrada (dose)
    \draw[->, very thick, verdeacido] (-1.5,1) -- (-0.1,1) node[midway, above, font=\tiny] {Dose};

    % Saida (eliminacao)
    \draw[->, very thick, red] (3.1,1) -- (4.5,1) node[midway, above, font=\tiny] {$k_e$};

\end{tikzpicture}
\end{center}

\column{0.5\textwidth}
\textbf{EDO:}
\destaque{
\begin{equation*}
\derivada{C}{t} = -k_e C
\end{equation*}
}

\textbf{Solução (dose IV):}
\begin{equation*}
C(t) = C_0 e^{-k_e t}
\end{equation*}

\textbf{Parâmetros:}
\begin{align*}
t_{1/2} &= \frac{\ln 2}{k_e} = \frac{0.693}{k_e}\\
\text{CL} &= k_e \cdot V_d\\
\text{AUC} &= \frac{C_0}{k_e} = \frac{\text{Dose}}{CL}
\end{align*}
\end{columns}
\end{frame}

\begin{frame}{Modelo de Dois Compartimentos}
\begin{center}
\begin{tikzpicture}[scale=0.8]
    % Compartimento central
    \draw[thick, azulescuro, fill=azulclaro!30] (0,0) rectangle (2.5,2);
    \node at (1.25,1.3) {\textbf{Central}};
    \node at (1.25,0.7) {$V_1$, $C_1(t)$};

    % Compartimento periferico
    \draw[thick, azulescuro, fill=verdeacido!20] (5,0) rectangle (7.5,2);
    \node at (6.25,1.3) {\textbf{Periférico}};
    \node at (6.25,0.7) {$V_2$, $C_2(t)$};

    % Dose
    \draw[->, very thick, verdeacido] (-1.5,1) -- (-0.1,1) node[midway, above, font=\tiny] {Dose};

    % Distribuicao
    \draw[->, thick, roxodna] (2.6,1.3) -- (4.9,1.3) node[midway, above, font=\tiny] {$k_{12}$};
    \draw[->, thick, roxodna] (4.9,0.7) -- (2.6,0.7) node[midway, below, font=\tiny] {$k_{21}$};

    % Eliminacao
    \draw[->, very thick, red] (1.25,-0.2) -- (1.25,-1.5) node[midway, right, font=\tiny] {$k_e$};

\end{tikzpicture}
\end{center}

\textbf{Sistema de EDOs:}
\destaque{
\begin{align*}
\derivada{C_1}{t} &= -(k_e + k_{12})C_1 + k_{21}C_2\\
\derivada{C_2}{t} &= k_{12}C_1 - k_{21}C_2
\end{align*}
}
\end{frame}

\begin{frame}[fragile]{Simulação PK: Modelo de Um Compartimento}
\begin{lstlisting}[language=Python, caption=Farmacocinética monocompartimental, basicstyle=\ttfamily\footnotesize]
import numpy as np
import matplotlib.pyplot as plt
from scipy.integrate import odeint

def one_compartment_pk(C, t, ke):
    """Modelo PK de 1 compartimento"""
    dCdt = -ke * C
    return dCdt

# Parametros
ke = 0.2      # constante de eliminacao (1/h)
C0 = 100      # concentracao inicial (mg/L)
t_half = 0.693 / ke  # meia-vida

# Tempo
t = np.linspace(0, 24, 200)

# Resolver
C = odeint(one_compartment_pk, C0, t, args=(ke,))

# Plotar
plt.figure(figsize=(10, 6))
plt.plot(t, C, 'b-', linewidth=2, label='[Farmaco]')
plt.axhline(y=C0/2, color='r', linestyle='--',
            label=f't_1/2 = {t_half:.2f} h')
plt.xlabel('Tempo (horas)')
plt.ylabel('Concentracao (mg/L)')
plt.title('Farmacocinetica: Modelo 1-Compartimento')
plt.legend()
plt.grid(True)
\end{lstlisting}
\end{frame}

\begin{frame}{Ensaios Clínicos e Biologia de Sistemas}
\begin{block}{Fases de Ensaios Clínicos}
\begin{itemize}
    \item \textbf{Fase I:} Segurança, farmacocinética (20-100 voluntários)
    \item \textbf{Fase II:} Eficácia preliminar, dose (100-500 pacientes)
    \item \textbf{Fase III:} Eficácia, efeitos adversos (1000-5000 pacientes)
    \item \textbf{Fase IV:} Monitoramento pós-comercialização
\end{itemize}
\end{block}

\vspace{0.3cm}

\begin{exampleblock}{Aplicações de Biologia de Sistemas}
\begin{itemize}
    \item \textbf{Design de ensaios:} Estratificação de pacientes por perfil molecular
    \item \textbf{Biomarcadores:} Identificação de marcadores de resposta/toxicidade
    \item \textbf{Ensaios \textit{in silico}:} "Digital twins" para prever resultados
    \item \textbf{Adaptive trials:} Ajustes baseados em dados intermediários
    \item \textbf{Análise de subgrupos:} Identificar respondedores vs não-respondedores
\end{itemize}
\end{exampleblock}
\end{frame}

% ============================================================================
% SEÇÃO 3: DOENÇAS COMPLEXAS
% ============================================================================
\section{Doenças Complexas}

\begin{frame}{Câncer como Doença Sistêmica}
\definicao{Visão Sistêmica do Câncer}{
Câncer não é apenas uma doença de um gene mutado, mas resulta de alterações em redes regulatórias, vias de sinalização e interações com o microambiente tumoral.
}

\vspace{0.3cm}

\begin{columns}
\column{0.5\textwidth}
\textbf{Características:}
\begin{itemize}
    \item Heterogeneidade genômica
    \item Plasticidade celular
    \item Evolução clonal
    \item Interações com estroma
    \item Resposta imune
    \item Metabolismo reprogramado
\end{itemize}

\column{0.5\textwidth}
\textbf{Níveis de Complexidade:}
\begin{itemize}
    \item Genético: mutações driver/passenger
    \item Epigenético: metilação, histonas
    \item Transcriptômico: expressão gênica
    \item Proteômico: sinalização
    \item Metabólico: Warburg, glutaminólise
    \item Microambiente: células imunes, CAFs
\end{itemize}
\end{columns}

\vspace{0.3cm}

\importante{
Abordagens de alvo único falham devido à redundância e plasticidade das redes tumorais
}
\end{frame}

\begin{frame}{Hallmarks of Cancer}
\begin{block}{10 Características Essenciais do Câncer (Hanahan \& Weinberg)}
\begin{columns}
\column{0.5\textwidth}
\textbf{Características Principais:}
\begin{enumerate}
    \item Sinais de crescimento sustentado
    \item Evasão de supressores de crescimento
    \item Resistência à morte celular
    \item Imortalidade replicativa
    \item Angiogênese induzida
    \item Invasão e metástase
\end{enumerate}

\column{0.5\textwidth}
\textbf{Características Emergentes:}
\begin{enumerate}
    \setcounter{enumi}{6}
    \item Reprogramação metabólica
    \item Evasão imune
\end{enumerate}

\vspace{0.3cm}

\textbf{Características Facilitadoras:}
\begin{enumerate}
    \setcounter{enumi}{8}
    \item Instabilidade genômica
    \item Inflamação promotora de tumor
\end{enumerate}
\end{columns}
\end{block}

\vspace{0.3cm}

\begin{exampleblock}{Implicações Terapêuticas}
Cada hallmark representa um alvo terapêutico potencial
\end{exampleblock}
\end{frame}

\begin{frame}{Redes Oncogênicas e Supressoras de Tumor}
\begin{center}
\begin{tikzpicture}[scale=0.9]
    % Oncogenes
    \node[draw, circle, fill=red!30, minimum size=0.7cm] (ras) at (0,2) {RAS};
    \node[draw, circle, fill=red!30, minimum size=0.7cm] (myc) at (2,2) {MYC};
    \node[draw, circle, fill=red!30, minimum size=0.7cm] (pi3k) at (1,3) {PI3K};

    % Supressores
    \node[draw, circle, fill=verdeacido!30, minimum size=0.7cm] (p53) at (0,0) {p53};
    \node[draw, circle, fill=verdeacido!30, minimum size=0.7cm] (rb) at (2,0) {RB};
    \node[draw, circle, fill=verdeacido!30, minimum size=0.7cm] (pten) at (1,-1) {PTEN};

    % Processos celulares
    \node[draw, rectangle, fill=azulclaro!30, minimum width=1.5cm] (prolif) at (5,2) {Proliferação};
    \node[draw, rectangle, fill=laranjacel!30, minimum width=1.5cm] (apop) at (5,0) {Apoptose};

    % Conexoes oncogenes
    \draw[->, thick, red] (ras) -- (prolif);
    \draw[->, thick, red] (myc) -- (prolif);
    \draw[->, thick, red] (pi3k) -- (prolif);
    \draw[-|, thick, red] (pi3k) to[bend left=20] (apop);

    % Conexoes supressores
    \draw[-|, thick, verdeacido] (p53) -- (prolif);
    \draw[-|, thick, verdeacido] (rb) -- (prolif);
    \draw[->, thick, verdeacido] (p53) -- (apop);
    \draw[->, thick, verdeacido] (pten) -- (apop);

    % Regulacao cruzada
    \draw[-|, thick, azulescuro, dashed] (pten) to[bend right=30] (pi3k);
    \draw[-|, thick, azulescuro, dashed] (p53) to[bend left=40] (myc);

    % Legenda
    \node[below, font=\tiny, text width=3cm, align=center] at (1,-2) {
        \textcolor{red}{Oncogenes ativam proliferação}\\
        \textcolor{verdeacido}{Supressores inibem proliferação}
    };
\end{tikzpicture}
\end{center}

\importante{
Câncer emerge do desequilíbrio entre oncogenes ativados e supressores inativados
}
\end{frame}

\begin{frame}{Metabolismo do Câncer: Efeito Warburg}
\begin{columns}
\column{0.5\textwidth}
\definicao{Efeito Warburg}{
Células cancerosas preferem glicólise aeróbica (menos eficiente) em vez de fosforilação oxidativa, mesmo na presença de oxigênio.
}

\vspace{0.3cm}

\textbf{Vantagens para tumor:}
\begin{itemize}
    \item Geração rápida de ATP
    \item Biomassa (nucleotídeos, lipídios)
    \item Acidificação do microambiente
    \item Evasão de apoptose mitocondrial
\end{itemize}

\column{0.5\textwidth}
\begin{block}{Reprogramação Metabólica}
\textbf{Vias alteradas:}
\begin{itemize}
    \item $\uparrow$ Captação de glicose (GLUT1)
    \item $\uparrow$ Glicólise (HK2, PKM2)
    \item $\uparrow$ Glutaminólise
    \item $\uparrow$ Síntese de lipídios (FASN)
    \item $\downarrow$ OXPHOS
\end{itemize}
\end{block}

\begin{exampleblock}{Alvos Metabólicos}
\begin{itemize}
    \item Inibidores de hexoquinase
    \item Inibidores de LDH-A
    \item Inibidores de glutaminase
\end{itemize}
\end{exampleblock}
\end{columns}
\end{frame}

\begin{frame}{Heterogeneidade Tumoral}
\begin{block}{Tipos de Heterogeneidade}
\begin{itemize}
    \item \textbf{Intratumoral:} diferentes clones no mesmo tumor
    \item \textbf{Intertumoral:} variação entre pacientes
    \item \textbf{Temporal:} evolução durante tratamento
    \item \textbf{Espacial:} gradientes no microambiente
\end{itemize}
\end{block}

\vspace{0.3cm}

\begin{center}
\begin{tikzpicture}[scale=0.8]
    % Tumor primario
    \draw[thick, fill=red!20] (0,0) circle (1.5);
    \node at (0,0) {\textbf{Tumor}};

    % Clones
    \fill[red!60] (-0.5,0.3) circle (0.3);
    \fill[red!80] (0.3,0.5) circle (0.25);
    \fill[red!40] (0.5,-0.4) circle (0.35);
    \fill[red!70] (-0.3,-0.5) circle (0.2);

    \node[font=\tiny] at (-0.5,0.3) {C1};
    \node[font=\tiny] at (0.3,0.5) {C2};
    \node[font=\tiny] at (0.5,-0.4) {C3};
    \node[font=\tiny] at (-0.3,-0.5) {C4};

    % Metastases
    \draw[->, thick] (1.3,0.8) -- (3,1.5);
    \draw[->, thick] (1.3,-0.8) -- (3,-1.5);

    \draw[thick, fill=red!60] (3.5,1.5) circle (0.5);
    \node[font=\tiny] at (3.5,1.5) {Met 1};

    \draw[thick, fill=red!40] (3.5,-1.5) circle (0.5);
    \node[font=\tiny] at (3.5,-1.5) {Met 2};

    % Legenda
    \node[below, font=\tiny, text width=4cm] at (1.5,-2.5) {
        Diferentes clones (C1-C4) têm perfis genômicos distintos\\
        Metástases derivam de clones específicos
    };
\end{tikzpicture}
\end{center}

\importante{
Heterogeneidade $\rightarrow$ resistência terapêutica $\rightarrow$ necessidade de terapias combinadas
}
\end{frame}

\begin{frame}{Doenças Cardiovasculares como Sistemas}
\begin{block}{Redes em Doenças Cardiovasculares}
Doenças cardiovasculares envolvem interações complexas entre:
\begin{itemize}
    \item Fatores genéticos (polimorfismos, mutações)
    \item Metabolismo lipídico (colesterol, triglicerídeos)
    \item Inflamação (citocinas, células imunes)
    \item Função endotelial (NO, vasodilatação)
    \item Coagulação (plaquetas, fatores de coagulação)
    \item Pressão arterial (RAAS, SNS)
\end{itemize}
\end{block}

\vspace{0.3cm}

\begin{columns}
\column{0.5\textwidth}
\textbf{Aterosclerose:}
\begin{itemize}
    \item Acúmulo de LDL oxidado
    \item Recrutamento de macrófagos
    \item Formação de células foam
    \item Placa aterosclerótica
    \item Inflamação crônica
\end{itemize}

\column{0.5\textwidth}
\textbf{Abordagens Sistêmicas:}
\begin{itemize}
    \item Identificação de módulos de risco
    \item Biomarcadores multi-ômicos
    \item Scores de risco poligênico
    \item Terapias multi-alvo
    \item Medicina preventiva personalizada
\end{itemize}
\end{columns}
\end{frame}

\begin{frame}{Diabetes e Síndrome Metabólica}
\begin{block}{Diabetes Tipo 2: Doença Sistêmica}
\textbf{Componentes interconectados:}
\begin{itemize}
    \item Resistência à insulina (músculo, fígado, tecido adiposo)
    \item Disfunção de células $\beta$ pancreáticas
    \item Metabolismo de glicose alterado
    \item Lipotoxicidade
    \item Inflamação crônica de baixo grau
    \item Estresse oxidativo
\end{itemize}
\end{block}

\vspace{0.3cm}

\begin{columns}
\column{0.5\textwidth}
\textbf{Síndrome Metabólica:}
\begin{enumerate}
    \item Obesidade abdominal
    \item Hiperglicemia
    \item Hipertensão
    \item Dislipidemia
    \item Estado pró-inflamatório
\end{enumerate}

\column{0.5\textwidth}
\begin{exampleblock}{Modelagem Sistêmica}
\begin{itemize}
    \item Modelos de homeostase glicose-insulina
    \item Redes de sinalização insulínica
    \item Integração multi-órgão
    \item Predição de progressão
    \item Resposta a intervenções
\end{itemize}
\end{exampleblock}
\end{columns}
\end{frame}

\begin{frame}{Doenças Neurodegenerativas}
\begin{columns}
\column{0.5\textwidth}
\begin{block}{Alzheimer}
\textbf{Características:}
\begin{itemize}
    \item Placas amiloides ($\beta$-amiloide)
    \item Emaranhados neurofibrilares (tau)
    \item Neuroinflamação
    \item Perda sináptica
    \item Disfunção mitocondrial
    \item Estresse oxidativo
\end{itemize}

\vspace{0.2cm}

\textbf{Fatores de risco:}
\begin{itemize}
    \item APOE4 (genético)
    \item Idade
    \item Fatores cardiovasculares
\end{itemize}
\end{block}

\column{0.5\textwidth}
\begin{block}{Parkinson}
\textbf{Características:}
\begin{itemize}
    \item Perda de neurônios dopaminérgicos
    \item Corpos de Lewy ($\alpha$-sinucleína)
    \item Disfunção mitocondrial
    \item Estresse do RE
    \item Neuroinflamação
\end{itemize}

\vspace{0.2cm}

\textbf{Genes associados:}
\begin{itemize}
    \item SNCA, LRRK2, PARK7
    \item GBA, PINK1, PARKIN
\end{itemize}
\end{block}
\end{columns}

\vspace{0.3cm}

\importante{
Doenças neurodegenerativas: falha em múltiplos processos celulares $\rightarrow$ cascata patológica
}
\end{frame}

\begin{frame}{Módulos de Doença (Disease Modules)}
\definicao{Módulo de Doença}{
Conjunto de genes/proteínas topologicamente próximos no interactoma humano cujas perturbações estão associadas a um fenótipo patológico comum.
}

\vspace{0.3cm}

\begin{center}
\begin{tikzpicture}[scale=0.8]
    % Rede global
    \foreach \i in {1,...,20} {
        \pgfmathsetmacro{\x}{rand*4}
        \pgfmathsetmacro{\y}{rand*4}
        \node[draw, circle, fill=azulclaro!20, inner sep=1pt, minimum size=0.2cm] (n\i) at (\x,\y) {};
    }

    % Modulo de doenca 1
    \draw[thick, red, dashed] (0.5,0.5) circle (1);
    \node[fill=red!30, circle, inner sep=2pt] at (0.5,0.5) {};
    \node[fill=red!30, circle, inner sep=2pt] at (0.8,1.2) {};
    \node[fill=red!30, circle, inner sep=2pt] at (1.3,0.7) {};
    \node[below, font=\tiny, text=red] at (0.5,-0.5) {Doença A};

    % Modulo de doenca 2
    \draw[thick, verdeacido, dashed] (3,3) circle (0.8);
    \node[fill=verdeacido!50, circle, inner sep=2pt] at (3,3) {};
    \node[fill=verdeacido!50, circle, inner sep=2pt] at (3.5,3.3) {};
    \node[fill=verdeacido!50, circle, inner sep=2pt] at (2.7,3.6) {};
    \node[above, font=\tiny, text=verdeacido] at (3,3.8) {Doença B};

    % Sobreposicao (comorbidade)
    \node[fill=yellow, circle, inner sep=2pt] at (2,2) {};
    \node[font=\tiny] at (2,1.5) {Comorbidade};

\end{tikzpicture}
\end{center}

\begin{exampleblock}{Aplicações}
\begin{itemize}
    \item Identificação de novos genes de doença
    \item Compreensão de comorbidades (sobreposição de módulos)
    \item Reposicionamento de fármacos
    \item Classificação molecular de doenças
\end{itemize}
\end{exampleblock}
\end{frame}

\begin{frame}{Classificação de Doenças Baseada em Redes}
\begin{block}{Abordagem Tradicional vs Sistêmica}
\begin{center}
\begin{tabular}{lcc}
\toprule
\textbf{Critério} & \textbf{Tradicional} & \textbf{Baseada em Rede} \\
\midrule
Base & Sintomas clínicos & Módulos moleculares \\
Foco & Órgão/tecido & Vias e redes \\
Resolução & Baixa & Alta (molecular) \\
Personalização & Limitada & Alta \\
Preditiva & Moderada & Alta \\
\bottomrule
\end{tabular}
\end{center}
\end{block}

\vspace{0.3cm}

\begin{exampleblock}{Exemplos}
\begin{itemize}
    \item \textbf{Câncer:} Classificação por perfil de expressão e mutações (The Cancer Genome Atlas - TCGA)
    \item \textbf{Diabetes:} Subtipos baseados em resistência à insulina vs disfunção de células $\beta$
    \item \textbf{Doenças inflamatórias:} Endotipos baseados em perfis de citocinas
\end{itemize}
\end{exampleblock}

\importante{
Doenças com apresentação clínica similar podem ter mecanismos moleculares distintos $\rightarrow$ requerem tratamentos diferentes
}
\end{frame}

% ============================================================================
% SEÇÃO 4: MEDICINA PERSONALIZADA
% ============================================================================
\section{Medicina Personalizada}

\begin{frame}{Farmacogenômica}
\definicao{Farmacogenômica}{
Estudo de como variações genéticas individuais influenciam a resposta a fármacos, incluindo eficácia e toxicidade.
}

\vspace{0.3cm}

\begin{block}{Fontes de Variabilidade na Resposta a Fármacos}
\begin{itemize}
    \item \textbf{Farmacocinética:} variantes em genes ADME (CYP450, transportadores)
    \item \textbf{Farmacodinâmica:} variantes em alvos terapêuticos (receptores, enzimas)
    \item \textbf{Doença:} heterogeneidade molecular da doença
    \item \textbf{Ambiente:} dieta, comorbidades, interações medicamentosas
\end{itemize}
\end{block}

\vspace{0.3cm}

\begin{columns}
\column{0.5\textwidth}
\textbf{Objetivos:}
\begin{itemize}
    \item Maximizar eficácia
    \item Minimizar toxicidade
    \item Otimizar dose
    \item Selecionar fármaco adequado
\end{itemize}

\column{0.5\textwidth}
\begin{exampleblock}{Benefícios}
\begin{itemize}
    \item Medicina personalizada
    \item Redução de eventos adversos
    \item Economia de custos
    \item Melhor adesão
\end{itemize}
\end{exampleblock}
\end{columns}
\end{frame}

\begin{frame}{Polimorfismos do Citocromo P450}
\begin{block}{Família CYP450}
Enzimas hepáticas responsáveis pelo metabolismo de ~75\% dos fármacos

\textbf{Principais isoformas:}
\begin{itemize}
    \item \textbf{CYP2D6:} ~25\% dos fármacos (antidepressivos, opioides, betabloqueadores)
    \item \textbf{CYP2C9:} warfarina, NSAIDs, hipoglicemiantes
    \item \textbf{CYP2C19:} clopidogrel, IBPs, antidepressivos
    \item \textbf{CYP3A4/5:} ~50\% dos fármacos (estatinas, imunossupressores)
\end{itemize}
\end{block}

\vspace{0.3cm}

\begin{block}{Fenótipos Metabólicos}
\begin{center}
\begin{tabular}{lll}
\toprule
\textbf{Fenótipo} & \textbf{Atividade} & \textbf{Implicação} \\
\midrule
Poor Metabolizer (PM) & Muito baixa & $\uparrow$ toxicidade, $\uparrow$ dose necessária (pró-fármacos) \\
Intermediate (IM) & Reduzida & Monitoramento \\
Normal (NM/EM) & Normal & Dose padrão \\
Ultra-rapid (UM) & Aumentada & $\downarrow$ eficácia, $\uparrow$ toxicidade (pró-fármacos) \\
\bottomrule
\end{tabular}
\end{center}
\end{block}
\end{frame}

\begin{frame}{Exemplo: Clopidogrel e CYP2C19}
\begin{columns}
\column{0.5\textwidth}
\textbf{Clopidogrel:}
\begin{itemize}
    \item Pró-fármaco antiagregante
    \item Previne trombose (pós-stent)
    \item Requer ativação por CYP2C19
\end{itemize}

\vspace{0.3cm}

\textbf{Polimorfismo CYP2C19:}
\begin{itemize}
    \item \textbf{*1/*1 (EM):} metabolizador normal
    \item \textbf{*2, *3 (PM):} perda de função
    \item \textbf{*17 (UM):} ganho de função
\end{itemize}

\vspace{0.3cm}

\importante{
PM têm $\uparrow$3x risco de eventos cardiovasculares
}

\column{0.5\textwidth}
\begin{center}
\begin{tikzpicture}[scale=0.7]
    % Eixos
    \draw[->] (0,0) -- (5,0) node[right, font=\tiny] {Tempo};
    \draw[->] (0,0) -- (0,4) node[above, font=\tiny] {[Metabólito ativo]};

    % Curvas
    \draw[thick, verdeacido] plot [smooth] coordinates {(0.5,0.2) (1,2.5) (2,3.5) (3,3.2) (4,2.5)};
    \draw[thick, red] plot [smooth] coordinates {(0.5,0.1) (1,0.8) (2,1.2) (3,1.1) (4,0.9)};
    \draw[thick, azulclaro] plot [smooth] coordinates {(0.5,0.15) (1,1.6) (2,2.3) (3,2.1) (4,1.7)};

    % Legenda
    \node[right, font=\tiny, text=verdeacido] at (4.2,3.2) {UM};
    \node[right, font=\tiny, text=azulclaro] at (4.2,2.1) {EM};
    \node[right, font=\tiny, text=red] at (4.2,1.1) {PM};

    % Zona terapeutica
    \draw[dashed, gray] (0,1.8) -- (5,1.8) node[right, font=\tiny] {Mín terapêutico};

\end{tikzpicture}
\end{center}

\vspace{0.3cm}

\textbf{Recomendação Clínica:}
\begin{itemize}
    \item PM: usar prasugrel ou ticagrelor (alternativas)
    \item EM: clopidogrel dose padrão
    \item UM: monitorar cuidadosamente
\end{itemize}
\end{columns}
\end{frame}

\begin{frame}{Descoberta de Biomarcadores}
\begin{block}{Tipos de Biomarcadores}
\begin{itemize}
    \item \textbf{Diagnóstico:} detectar presença de doença
    \item \textbf{Prognóstico:} prever evolução da doença
    \item \textbf{Preditivo:} prever resposta a terapia
    \item \textbf{Farmacodinâmico:} medir efeito do tratamento
    \item \textbf{Segurança:} detectar toxicidade
\end{itemize}
\end{block}

\vspace{0.3cm}

\begin{exampleblock}{Pipeline de Descoberta}
\begin{enumerate}
    \item \textbf{Discovery:} análise ômica em larga escala (casos vs controles)
    \item \textbf{Verificação:} validação em coorte independente
    \item \textbf{Qualificação:} ensaios prospectivos
    \item \textbf{Implementação:} testes clínicos validados
\end{enumerate}
\end{exampleblock}

\begin{columns}
\column{0.5\textwidth}
\textbf{Abordagens Sistêmicas:}
\begin{itemize}
    \item Assinaturas multi-gene
    \item Perfis de proteínas/metabólitos
    \item Integração multi-ômica
\end{itemize}

\column{0.5\textwidth}
\textbf{Desafios:}
\begin{itemize}
    \item Reprodutibilidade
    \item Validação independente
    \item Padronização de ensaios
    \item Custo-benefício
\end{itemize}
\end{columns}
\end{frame}

\begin{frame}{Estratificação de Pacientes}
\begin{block}{Conceito}
Classificar pacientes em subgrupos com base em características moleculares, clínicas ou demográficas para tratamento personalizado.
\end{block}

\vspace{0.3cm}

\begin{columns}
\column{0.5\textwidth}
\textbf{Métodos:}
\begin{itemize}
    \item Clustering (K-means, hierárquico)
    \item Análise de componentes principais (PCA)
    \item Non-negative matrix factorization (NMF)
    \item Machine learning supervisionado
    \item Network-based stratification
\end{itemize}

\vspace{0.3cm}

\textbf{Dados para estratificação:}
\begin{itemize}
    \item Genômicos (mutações, CNVs)
    \item Transcriptômicos (expressão gênica)
    \item Proteômicos
    \item Metabólicos
    \item Clínicos
\end{itemize}

\column{0.5\textwidth}
\begin{exampleblock}{Aplicações}
\begin{itemize}
    \item \textbf{Ensaios clínicos:} enriquecer para respondedores
    \item \textbf{Diagnóstico:} classificar subtipos
    \item \textbf{Prognóstico:} identificar alto vs baixo risco
    \item \textbf{Terapia:} selecionar tratamento ótimo
\end{itemize}
\end{exampleblock}

\vspace{0.3cm}

\importante{
Objetivo: maximizar benefício clínico, minimizar exposição desnecessária a toxicidade
}
\end{columns}
\end{frame}

\begin{frame}{Diagnósticos Baseados em Ômicas}
\begin{block}{Testes Multi-Ômicos Clínicos}
\textbf{Genômicos:}
\begin{itemize}
    \item Sequenciamento de painel de genes (câncer)
    \item Whole-exome sequencing (WES)
    \item Whole-genome sequencing (WGS)
    \item Liquid biopsy (ctDNA)
\end{itemize}

\textbf{Transcriptômicos:}
\begin{itemize}
    \item Oncotype DX (câncer de mama)
    \item MammaPrint (prognóstico de mama)
    \item Prosigna (PAM50)
\end{itemize}

\textbf{Proteômicos/Metabólicos:}
\begin{itemize}
    \item Painéis de citocinas
    \item Perfis metabólicos (diabetes, erros inatos)
\end{itemize}
\end{block}

\vspace{0.3cm}

\begin{exampleblock}{Vantagens}
Informação molecular detalhada $\rightarrow$ decisões terapêuticas precisas
\end{exampleblock}
\end{frame}

\begin{frame}{Subtipos de Câncer de Mama}
\begin{block}{Classificação Molecular (PAM50)}
Baseada em perfil de expressão de 50 genes:

\begin{center}
\begin{tabular}{llll}
\toprule
\textbf{Subtipo} & \textbf{Marcadores} & \textbf{Prognóstico} & \textbf{Tratamento} \\
\midrule
Luminal A & ER+, PR+, HER2-, Ki67 baixo & Melhor & Hormonioterapia \\
Luminal B & ER+, HER2+/-, Ki67 alto & Intermediário & Hormônio + quimio \\
HER2+ & ER-, PR-, HER2+ & Intermediário & Trastuzumab + quimio \\
Basal-like & ER-, PR-, HER2- (triplo neg.) & Pior & Quimioterapia \\
\bottomrule
\end{tabular}
\end{center}
\end{block}

\vspace{0.3cm}

\importante{
Mesmo órgão, histologia similar $\rightarrow$ diferentes doenças moleculares $\rightarrow$ tratamentos distintos
}

\begin{exampleblock}{Implicação Clínica}
Teste de expressão gênica determina:
\begin{itemize}
    \item Necessidade de quimioterapia adjuvante
    \item Escolha de terapia-alvo (trastuzumab para HER2+)
    \item Prognóstico e risco de recorrência
\end{itemize}
\end{exampleblock}
\end{frame}

\begin{frame}{Seleção de Tratamento por Perfil Molecular}
\begin{center}
\begin{tikzpicture}[scale=0.8]
    % Paciente
    \node[draw, circle, fill=azulclaro!30, minimum size=1.2cm] (patient) at (0,0) {\textbf{Paciente}};

    % Perfil molecular
    \node[draw, rectangle, fill=roxodna!20, text width=2cm, align=center] (profile) at (3,0) {Perfil\\Molecular};

    % Analise
    \node[draw, rectangle, fill=verdeacido!20, text width=2cm, align=center] (analysis) at (6,0) {Análise\\Sistêmica};

    % Tratamentos
    \node[draw, rectangle, fill=laranjacel!30, text width=1.5cm, align=center] (t1) at (9,1.5) {Terapia A};
    \node[draw, rectangle, fill=laranjacel!30, text width=1.5cm, align=center] (t2) at (9,0) {Terapia B};
    \node[draw, rectangle, fill=laranjacel!30, text width=1.5cm, align=center] (t3) at (9,-1.5) {Terapia C};

    % Conexoes
    \draw[->, thick] (patient) -- (profile);
    \draw[->, thick] (profile) -- (analysis);
    \draw[->, thick] (analysis) -- (t1) node[midway, above, font=\tiny] {HER2+};
    \draw[->, thick, red, dashed] (analysis) -- (t2);
    \draw[->, thick, red, dashed] (analysis) -- (t3);

    % Genomics, transcriptomics
    \node[below, font=\tiny, text width=2cm, align=center] at (3,-1) {Genômica\\Transcriptômica\\Proteômica};

    % Machine learning
    \node[above, font=\tiny] at (6,0.8) {ML/Redes};

\end{tikzpicture}
\end{center}

\vspace{0.3cm}

\begin{exampleblock}{Exemplos Clínicos}
\begin{itemize}
    \item \textbf{EGFR mutado:} erlotinib/gefitinib (câncer de pulmão)
    \item \textbf{BRAF V600E:} vemurafenib (melanoma)
    \item \textbf{ALK fusão:} crizotinib (câncer de pulmão)
    \item \textbf{BCR-ABL:} imatinib (leucemia mieloide crônica)
\end{itemize}
\end{exampleblock}
\end{frame}

\begin{frame}{Digital Twins e Ensaios Clínicos \textit{In Silico}}
\definicao{Digital Twin}{
Réplica computacional personalizada de um paciente individual, baseada em dados multi-ômicos, clínicos e fisiológicos, que permite simular respostas a intervenções terapêuticas.
}

\vspace{0.3cm}

\begin{columns}
\column{0.5\textwidth}
\textbf{Componentes:}
\begin{itemize}
    \item Dados do paciente (genoma, transcriptoma, metaboloma, clínicos)
    \item Modelos mecanísticos (redes, EDOs)
    \item Parâmetros personalizados
    \item Simulação computacional
\end{itemize}

\vspace{0.3cm}

\textbf{Aplicações:}
\begin{itemize}
    \item Predição de resposta a fármacos
    \item Otimização de dose
    \item Identificação de efeitos adversos
    \item Testes de terapias combinadas
\end{itemize}

\column{0.5\textwidth}
\begin{exampleblock}{Ensaios \textit{In Silico}}
\begin{itemize}
    \item Testar milhares de cenários rapidamente
    \item Reduzir necessidade de ensaios animais
    \item Acelerar desenvolvimento de fármacos
    \item Stratificação pré-clínica
\end{itemize}
\end{exampleblock}

\vspace{0.3cm}

\begin{alertblock}{Desafios}
\begin{itemize}
    \item Validação clínica
    \item Complexidade dos modelos
    \item Integração de dados
    \item Incerteza paramétrica
\end{itemize}
\end{alertblock}
\end{columns}
\end{frame}

\begin{frame}{Desafios em Medicina Personalizada}
\begin{alertblock}{Desafios Técnicos}
\begin{itemize}
    \item \textbf{Integração de dados:} multi-ômicos heterogêneos
    \item \textbf{Interpretação:} variantes de significado incerto (VUS)
    \item \textbf{Validação:} necessidade de estudos prospectivos
    \item \textbf{Padronização:} protocolos laboratoriais e análise
    \item \textbf{Custo:} sequenciamento e análise ainda caros
\end{itemize}
\end{alertblock}

\vspace{0.3cm}

\begin{alertblock}{Desafios Éticos e Sociais}
\begin{itemize}
    \item \textbf{Privacidade:} dados genômicos sensíveis
    \item \textbf{Consentimento:} uso de dados, achados incidentais
    \item \textbf{Equidade:} acesso desigual a tecnologias
    \item \textbf{Discriminação:} genética por seguradoras/empregadores
    \item \textbf{Educação:} médicos e pacientes precisam compreender resultados
\end{itemize}
\end{alertblock}

\importante{
Medicina personalizada requer framework regulatório, ético e de treinamento profissional adequado
}
\end{frame}

% ============================================================================
% SEÇÃO 5: APLICAÇÕES CLÍNICAS
% ============================================================================
\section{Aplicações Clínicas}

\begin{frame}{Estudo de Caso: Câncer de Mama HER2+ e Trastuzumab}
\begin{block}{Contexto Clínico}
\begin{itemize}
    \item 15-20\% dos cânceres de mama superexpressam HER2
    \item HER2 (ERBB2): receptor tirosina-quinase
    \item Superexpressão $\rightarrow$ proliferação descontrolada
    \item Prognóstico pior (antes de terapias-alvo)
\end{itemize}
\end{block}

\vspace{0.3cm}

\begin{columns}
\column{0.5\textwidth}
\textbf{Trastuzumab (Herceptin):}
\begin{itemize}
    \item Anticorpo monoclonal anti-HER2
    \item Bloqueia sinalização HER2
    \item Induz resposta imune (ADCC)
    \item Primeiro sucesso de medicina personalizada em oncologia
\end{itemize}

\vspace{0.3cm}

\textbf{Impacto:}
\begin{itemize}
    \item $\downarrow$ 50\% mortalidade
    \item Sobrevida global melhorada
    \item Usado em adjuvância e metastático
\end{itemize}

\column{0.5\textwidth}
\begin{center}
\begin{tikzpicture}[scale=0.7]
    % Membrana
    \draw[thick] (0,0) -- (4,0);

    % HER2
    \draw[thick, fill=roxodna!30] (1,0) rectangle (1.5,0.3);
    \draw[thick, fill=roxodna!30] (1.2,-0.3) -- (1.2,0);
    \draw[thick, fill=roxodna!30] (1.2,-0.3) circle (0.15);
    \node[below, font=\tiny] at (1.2,-0.5) {HER2};

    % Trastuzumab
    \draw[thick, verdeacido, ->] (1.2,1.5) -- (1.2,0.5);
    \node[draw, circle, fill=verdeacido!30] (tras) at (1.2,2) {mAb};
    \node[right, font=\tiny] at (1.7,2) {Trastuzumab};

    % Sinalizacao bloqueada
    \draw[thick, red, dashed] (1.2,-0.5) -- (1.2,-1.5);
    \node[below, font=\tiny] at (1.2,-1.7) {Bloqueio};
    \draw[-|, thick, red] (1.2,-1) -- (2.5,-1);

    % Proliferacao
    \node[right, font=\tiny, text width=1.5cm] at (2.5,-1.2) {$\downarrow$ Proliferação};

\end{tikzpicture}
\end{center}

\vspace{0.3cm}

\importante{
Teste HER2 é mandatório em câncer de mama $\rightarrow$ decisão terapêutica
}
\end{columns}
\end{frame}

\begin{frame}{Doenças Infecciosas: Abordagem Sistêmica}
\begin{block}{Interações Hospedeiro-Patógeno}
Visão sistêmica considera:
\begin{itemize}
    \item Redes moleculares do patógeno
    \item Resposta imune do hospedeiro
    \item Interações proteína-proteína (patógeno-hospedeiro)
    \item Modulação de vias do hospedeiro por patógeno
    \item Evolução e adaptação do patógeno
\end{itemize}
\end{block}

\vspace{0.3cm}

\begin{columns}
\column{0.5\textwidth}
\textbf{Aplicações:}
\begin{itemize}
    \item Identificação de alvos terapêuticos (essenciais para patógeno)
    \item Reposicionamento de fármacos
    \item Previsão de virulência
    \item Assinaturas de resposta imune
    \item Estratégias de vacinação
\end{itemize}

\column{0.5\textwidth}
\begin{exampleblock}{Exemplo: Tuberculose}
\begin{itemize}
    \item Rede metabólica de \textit{M. tuberculosis}
    \item Identificar enzimas essenciais
    \item Testar \textit{in silico} inibidores
    \item Prever resistência a fármacos
    \item Biomarcadores de progressão latente $\rightarrow$ ativa
\end{itemize}
\end{exampleblock}
\end{columns}
\end{frame}

\begin{frame}{Resistência a Antibióticos como Problema Sistêmico}
\begin{block}{Mecanismos de Resistência}
\begin{itemize}
    \item Mutações em alvos (RNA polimerase, ribossomo)
    \item Bombas de efluxo (superexpressão)
    \item Enzimas inativadoras ($\beta$-lactamases)
    \item Modificação de alvos (metilação de RNA ribossomal)
    \item Biofilmes (tolerância)
\end{itemize}
\end{block}

\vspace{0.3cm}

\begin{columns}
\column{0.5\textwidth}
\textbf{Abordagens Sistêmicas:}
\begin{itemize}
    \item Análise de redes de resistência
    \item Identificar vulnerabilidades colaterais
    \item Terapias combinadas racionais
    \item Predição de evolução de resistência
    \item Desenvolvimento de novos antibióticos
\end{itemize}

\column{0.5\textwidth}
\begin{alertblock}{Estratégias}
\begin{itemize}
    \item \textbf{Collateral sensitivity:} resistência a A causa sensibilidade a B
    \item \textbf{Terapias alternadas:} reduzir pressão seletiva
    \item \textbf{Alvos não-essenciais:} virulência sem morte bacteriana
    \item \textbf{Imunomodulação:} fortalecer resposta do hospedeiro
\end{itemize}
\end{alertblock}
\end{columns}

\importante{
Resistência antimicrobiana é uma das maiores ameaças à saúde global
}
\end{frame}

\begin{frame}{Imunoterapia e Imunologia de Sistemas}
\begin{block}{Imunoterapia do Câncer}
Estratégias que ativam o sistema imune para eliminar células tumorais:
\begin{itemize}
    \item \textbf{Checkpoint inhibitors:} anti-PD-1, anti-CTLA-4 (bloqueiam sinais inibitórios)
    \item \textbf{CAR-T cells:} células T geneticamente modificadas
    \item \textbf{Vacinas terapêuticas:} antígenos tumorais
    \item \textbf{Oncolytic viruses:} vírus que infectam e lisam tumor
\end{itemize}
\end{block}

\vspace{0.3cm}

\begin{columns}
\column{0.5\textwidth}
\textbf{Desafios:}
\begin{itemize}
    \item Apenas 20-40\% respondem
    \item Efeitos adversos (autoimunidade)
    \item Microambiente imunossupressor
    \item Heterogeneidade tumoral
\end{itemize}

\column{0.5\textwidth}
\textbf{Imunologia de Sistemas:}
\begin{itemize}
    \item Perfil de infiltrado imune (hot vs cold tumors)
    \item Assinaturas de resposta
    \item Carga mutacional tumoral (TMB)
    \item Expressão de PD-L1
    \item Modelagem de resposta imune
\end{itemize}
\end{columns}

\vspace{0.3cm}

\importante{
Biomarcadores preditivos são essenciais para selecionar pacientes que se beneficiarão
}
\end{frame}

\begin{frame}{Terapia com Células CAR-T}
\begin{block}{Conceito}
Células T do paciente são geneticamente modificadas para expressar receptor de antígeno quimérico (CAR) que reconhece antígeno tumoral específico.
\end{block}

\vspace{0.3cm}

\begin{center}
\begin{tikzpicture}[scale=0.75]
    % Etapa 1: coleta
    \node[draw, circle, fill=azulclaro!30] (patient) at (0,0) {Paciente};
    \draw[->, thick] (1,0) -- (2,0);
    \node[draw, rectangle, fill=roxodna!20] (tcell) at (3,0) {Células T};

    % Etapa 2: modificacao
    \draw[->, thick] (4,0) -- (5,0);
    \node[draw, rectangle, fill=verdeacido!20, text width=1.5cm, align=center] (mod) at (6.5,0) {Modificação\\Genética\\(CAR)};

    % Etapa 3: expansao
    \draw[->, thick] (8,0) -- (9,0);
    \node[draw, rectangle, fill=laranjacel!20, text width=1.5cm, align=center] (exp) at (10.5,0) {Expansão\\\textit{ex vivo}};

    % Etapa 4: reinfusao
    \draw[->, thick] (12,0) -- (13,0);
    \node[draw, circle, fill=azulclaro!30] (back) at (14,0) {Paciente};

    % Resultado: tumor
    \node[below=0.5cm of back, draw, rectangle, fill=red!30] (tumor) {Tumor};
    \draw[->, thick, red] (back) -- (tumor) node[midway, right, font=\tiny] {CAR-T ataca};

\end{tikzpicture}
\end{center}

\vspace{0.3cm}

\begin{columns}
\column{0.5\textwidth}
\textbf{Sucessos:}
\begin{itemize}
    \item Leucemias/linfomas B (CD19)
    \item Taxa de remissão completa >80\%
    \item FDA aprovado (Kymriah, Yescarta)
\end{itemize}

\column{0.5\textwidth}
\textbf{Modelagem Sistêmica:}
\begin{itemize}
    \item Dinâmica CAR-T vs tumor
    \item Predição de síndrome de liberação de citocinas (CRS)
    \item Otimização de dose
    \item Design de CARs de próxima geração
\end{itemize}
\end{columns}
\end{frame}

\begin{frame}{Design de Vacinas: Abordagem Sistêmica}
\begin{block}{Vacinologia de Sistemas}
Uso de abordagens ômicas e modelagem para:
\begin{itemize}
    \item Identificar antígenos imunogênicos e protetores
    \item Prever resposta imune a candidatos vacinais
    \item Otimizar formulações e adjuvantes
    \item Estratificar respondedores vs não-respondedores
    \item Compreender correlatos de proteção
\end{itemize}
\end{block}

\vspace{0.3cm}

\begin{columns}
\column{0.5\textwidth}
\textbf{Pipeline:}
\begin{enumerate}
    \item Genômica de patógeno
    \item Predição de epítopos (MHC binding)
    \item Análise de conservação
    \item Teste \textit{in silico} de imunogenicidade
    \item Validação experimental
    \item Ensaios clínicos
\end{enumerate}

\column{0.5\textwidth}
\begin{exampleblock}{Exemplo: Vacinas mRNA (COVID-19)}
\begin{itemize}
    \item Análise genômica rápida do SARS-CoV-2
    \item Identificação da proteína Spike (S)
    \item Design de mRNA otimizado
    \item Predição de resposta imune
    \item Desenvolvimento em tempo recorde (<1 ano)
\end{itemize}
\end{exampleblock}
\end{columns}

\importante{
Biologia de sistemas acelerou desenvolvimento de vacinas COVID-19
}
\end{frame}

\begin{frame}{COVID-19 e Resposta de Biologia de Sistemas}
\begin{block}{Contribuições de Systems Biology na Pandemia}
\begin{itemize}
    \item \textbf{Genômica viral:} sequenciamento rápido, rastreamento de variantes
    \item \textbf{Reposicionamento de fármacos:} análise de redes identificou candidatos (remdesivir, dexametasona)
    \item \textbf{Resposta imune:} perfis de citocinas, assinaturas de gravidade (tempestade de citocinas)
    \item \textbf{Interações proteína-proteína:} mapa de interações SARS-CoV-2-humano
    \item \textbf{Modelagem epidemiológica:} previsão de disseminação, impacto de intervenções
    \item \textbf{Vacinas:} design racional (mRNA, vetores virais)
\end{itemize}
\end{block}

\vspace{0.3cm}

\begin{exampleblock}{Biomarcadores de Gravidade}
\begin{itemize}
    \item IL-6, D-dímero, ferritina $\uparrow$ em casos graves
    \item Linfopenia, neutrofilia
    \item Perfis transcriptômicos distinguem progressão
\end{itemize}
\end{exampleblock}
\end{frame}

\begin{frame}[fragile]{Exemplo: Estratificação de Pacientes com Clustering}
\begin{lstlisting}[language=Python, caption=Análise de clusters em dados clínicos, basicstyle=\ttfamily\footnotesize]
import numpy as np
import matplotlib.pyplot as plt
from sklearn.cluster import KMeans
from sklearn.preprocessing import StandardScaler
from sklearn.decomposition import PCA

# Dados simulados: pacientes x features (omicas + clinicas)
np.random.seed(42)
n_patients = 100
data = np.random.randn(n_patients, 10)  # 10 features

# Normalizar
scaler = StandardScaler()
data_scaled = scaler.fit_transform(data)

# Clustering K-means
kmeans = KMeans(n_clusters=3, random_state=42)
clusters = kmeans.fit_predict(data_scaled)

# PCA para visualizacao
pca = PCA(n_components=2)
data_pca = pca.transform(data_scaled)

# Plotar
plt.figure(figsize=(8, 6))
scatter = plt.scatter(data_pca[:, 0], data_pca[:, 1],
                       c=clusters, cmap='viridis', s=50)
plt.xlabel('PC1')
plt.ylabel('PC2')
plt.title('Estratificacao de Pacientes (K-means)')
plt.colorbar(scatter, label='Cluster')
plt.legend(['Subgrupo 1', 'Subgrupo 2', 'Subgrupo 3'])
\end{lstlisting}
\end{frame}

% ============================================================================
% SEÇÃO 6: EXERCÍCIOS
% ============================================================================
\section{Exercícios}

\begin{frame}{Exercícios - Parte 1: Farmacocinética}
\begin{block}{Questão 1: Modelo PK}
Um paciente recebe dose IV bolus de 500 mg de um fármaco. O volume de distribuição é $V_d = 50$ L e a meia-vida é $t_{1/2} = 4$ h.

\begin{enumerate}
    \item Calcule a constante de eliminação $k_e$
    \item Calcule o clearance (CL)
    \item Qual será a concentração após 8 horas?
    \item Calcule a AUC (área sob a curva)
\end{enumerate}
\end{block}

\begin{block}{Questão 2: Modelo de Dois Compartimentos}
Implemente um modelo farmacocinético de dois compartimentos em Python e simule a concentração ao longo de 24 horas. Use os seguintes parâmetros:
\begin{itemize}
    \item $k_e = 0.15$ h$^{-1}$, $k_{12} = 0.3$ h$^{-1}$, $k_{21} = 0.2$ h$^{-1}$
    \item Dose = 1000 mg IV, $C_1(0) = 20$ mg/L, $C_2(0) = 0$
\end{itemize}
\end{block}
\end{frame}

\begin{frame}{Exercícios - Parte 2: Redes de Doenças}
\begin{block}{Questão 3: Análise de Rede DTI}
Dado o seguinte grafo bipartido de interações fármaco-alvo:
\begin{itemize}
    \item Fármaco A: alvos T1, T2, T3
    \item Fármaco B: alvos T2, T4
    \item Fármaco C: alvos T3, T4, T5
\end{itemize}

\begin{enumerate}
    \item Construa a matriz de adjacência
    \item Calcule o grau de promiscuidade de cada fármaco
    \item Identifique alvos compartilhados (potencial para efeitos adversos similares)
    \item Sugira reposicionamento de fármacos baseado em sobreposição de alvos
\end{enumerate}
\end{block}

\begin{block}{Questão 4: Módulos de Doença}
Use algoritmo de Louvain ou Girvan-Newman para detectar módulos em uma rede de interações proteína-proteína associada a doença cardiovascular. Interprete os módulos funcionalmente.
\end{block}
\end{frame}

\begin{frame}{Exercícios - Parte 3: Medicina Personalizada}
\begin{block}{Questão 5: Farmacogenômica}
Um paciente é metabolizador pobre (PM) de CYP2D6 e recebe codeína (pró-fármaco convertido em morfina por CYP2D6).

\begin{enumerate}
    \item Preveja o efeito no manejo da dor
    \item Qual ajuste de dose ou medicação alternativa você recomendaria?
    \item Como você explicaria ao paciente a importância do teste farmacogenômico?
\end{enumerate}
\end{block}

\begin{block}{Questão 6: Estratificação de Pacientes}
Implemente clustering K-means em um dataset de expressão gênica de câncer de mama (use dataset público como TCGA ou simulado). Identifique subgrupos moleculares e compare com subtipos clínicos conhecidos (Luminal A/B, HER2+, Basal).
\end{block}
\end{frame}

\begin{frame}{Projeto Prático: Análise de Dados Clínicos Ômicos}
\begin{alertblock}{Projeto Integrador}
\textbf{Objetivo:} Analisar dados multi-ômicos de pacientes com câncer para identificar biomarcadores e estratificar pacientes.

\textbf{Etapas:}
\begin{enumerate}
    \item \textbf{Dados:} Baixar dados de expressão gênica e clínicos (TCGA, GEO, cBioPortal)
    \item \textbf{Pré-processamento:} Normalização, controle de qualidade, filtro de variância
    \item \textbf{Análise exploratória:} PCA, heatmaps, clustering hierárquico
    \item \textbf{Clustering:} K-means ou NMF para identificar subgrupos
    \item \textbf{Biomarcadores:} Análise diferencial (DESeq2, limma) entre subgrupos
    \item \textbf{Sobrevida:} Curvas de Kaplan-Meier, teste log-rank
    \item \textbf{Validação:} Dataset independente
    \item \textbf{Interpretação biológica:} Enriquecimento funcional (GO, KEGG)
\end{enumerate}

\textbf{Entregável:} Relatório com código Python/R, figuras e interpretação clínica
\end{alertblock}
\end{frame}

% ============================================================================
% SEÇÃO 7: REFERÊNCIAS
% ============================================================================
\section{Referências}

\begin{frame}{Referências Principais}
\begin{thebibliography}{99}
\bibitem{barabasi2011} Barabási A.L., Gulbahce N., Loscalzo J. (2011). \textit{Network medicine: a network-based approach to human disease}. Nature Reviews Genetics 12:56-68.

\bibitem{auffray2009} Auffray C., Chen Z., Hood L. (2009). \textit{Systems medicine: the future of medical genomics and healthcare}. Genome Medicine 1:2.

\bibitem{schadt2009} Schadt E.E. (2009). \textit{Molecular networks as sensors and drivers of common human diseases}. Nature 461:218-223.

\bibitem{zhou2014} Zhou X., Menche J., Barabási A.L., Sharma A. (2014). \textit{Human symptoms-disease network}. Nature Communications 5:4212.

\bibitem{hanahan2011} Hanahan D., Weinberg R.A. (2011). \textit{Hallmarks of cancer: the next generation}. Cell 144:646-674.
\end{thebibliography}
\end{frame}

\begin{frame}{Referências: Network Medicine e Farmacologia}
\begin{thebibliography}{99}
\bibitem{hopkins2008} Hopkins A.L. (2008). \textit{Network pharmacology: the next paradigm in drug discovery}. Nature Chemical Biology 4:682-690.

\bibitem{zhao2012} Zhao S., Iyengar R. (2012). \textit{Systems pharmacology: network analysis to identify multiscale mechanisms of drug action}. Annual Review of Pharmacology and Toxicology 52:505-521.

\bibitem{csermely2013} Csermely P., Korcsmáros T., Kiss H.J., et al. (2013). \textit{Structure and dynamics of molecular networks: a novel paradigm of drug discovery}. Pharmacology \& Therapeutics 138:333-408.

\bibitem{menche2015} Menche J., Sharma A., Kitsak M., et al. (2015). \textit{Disease networks: uncovering disease-disease relationships through the incomplete interactome}. Science 347:1257601.
\end{thebibliography}
\end{frame}

\begin{frame}{Referências: Medicina Personalizada e Ômicas}
\begin{thebibliography}{99}
\bibitem{ashley2016} Ashley E.A. (2016). \textit{Towards precision medicine}. Nature Reviews Genetics 17:507-522.

\bibitem{roden2015} Roden D.M., McLeod H.L., Relling M.V., et al. (2019). \textit{Pharmacogenomics}. The Lancet 394:521-532.

\bibitem{chen2012} Chen R., Snyder M. (2013). \textit{Promise of personalized omics to precision medicine}. Wiley Interdisciplinary Reviews: Systems Biology and Medicine 5:73-82.

\bibitem{brunet2004} Brunet J.P., et al. (2004). \textit{Metagenes and molecular pattern discovery using matrix factorization}. PNAS 101:4164-4169.
\end{thebibliography}
\end{frame}

\begin{frame}{Bases de Dados e Recursos}
\begin{block}{Bases de Dados Clínicas e Farmacogenômicas}
\begin{itemize}
    \item \textbf{DrugBank:} \url{https://go.drugbank.com} --- fármacos, alvos, PK/PD
    \item \textbf{TCGA:} \url{https://portal.gdc.cancer.gov} --- dados multi-ômicos de câncer
    \item \textbf{UK Biobank:} \url{https://www.ukbiobank.ac.uk} --- dados de saúde populacional
    \item \textbf{ClinVar:} \url{https://www.ncbi.nlm.nih.gov/clinvar} --- variantes clínicas
    \item \textbf{PharmGKB:} \url{https://www.pharmgkb.org} --- farmacogenômica
\end{itemize}
\end{block}

\begin{block}{Ferramentas de Análise}
\begin{itemize}
    \item \textbf{cBioPortal:} \url{https://www.cbioportal.org} --- análise de dados de câncer
    \item \textbf{STRING:} \url{https://string-db.org} --- interações proteína-proteína
    \item \textbf{Cytoscape:} \url{https://cytoscape.org} --- visualização de redes
    \item \textbf{GSEA:} \url{https://www.gsea-msigdb.org} --- enriquecimento funcional
\end{itemize}
\end{block}
\end{frame}

\begin{frame}{Leitura Complementar}
\begin{block}{Livros e Reviews}
\begin{itemize}
    \item Loscalzo J., Barabási A.L., Silverman E.K. (2017). \textit{Network Medicine: Complex Systems in Human Disease and Therapeutics}. Harvard University Press.
    \item Zhu H., Xia M. (2016). \textit{Advances in Systems Biology}. Springer.
    \item Wolkenhauer O. (2013). \textit{Systems Medicine}. Springer.
\end{itemize}
\end{block}

\begin{block}{Artigos de Revisão Importantes}
\begin{itemize}
    \item Hood L., Friend S.H. (2011). \textit{Predictive, personalized, preventive, participatory (P4) cancer medicine}. Nature Reviews Clinical Oncology 8:184-187.
    \item Ideker T., Krogan N.J. (2012). \textit{Differential network biology}. Molecular Systems Biology 8:565.
    \item Califano A., Alvarez M.J. (2017). \textit{The recurrent architecture of tumour initiation, progression and drug sensitivity}. Nature Reviews Cancer 17:116-130.
\end{itemize}
\end{block}
\end{frame}

% ============================================================================
% SLIDE FINAL
% ============================================================================
\begin{frame}
\begin{center}
{\Huge Obrigado!}

\vspace{1cm}

{\Large Capítulo 13: Systems Biology in Medicine}

\vspace{1cm}

\textbf{Tópicos Abordados:}\\
Descoberta de fármacos $\cdot$ Doenças complexas\\
Medicina personalizada $\cdot$ Aplicações clínicas

\vspace{1cm}

\textit{Dúvidas? Perguntas?}

\vspace{0.5cm}

\textit{A medicina do futuro é sistêmica, preditiva, personalizada e participativa}
\end{center}
\end{frame}

\end{document}
